\documentclass[12pt]{article}


\usepackage[pdftex,pdfpagelabels,bookmarks,hyperindex,hyperfigures]{hyperref}


\usepackage[margin=1in]{geometry}
\usepackage{float}
\usepackage{adjustbox}
\usepackage[table]{xcolor}
\usepackage{cite}
\usepackage{amsmath,amssymb,amsfonts}
\usepackage{algorithmic}
\usepackage{graphicx}
\usepackage{textcomp}
\usepackage{xcolor}
\usepackage{pgfplots}
\usepackage{fancyhdr}
\usepackage{color}
\usepackage{listings}

\setlength{\parskip}{1em} % Add spacing between paragraphs
\setlength{\parindent}{0em} % Remove indentation at the start of paragraphs


\pagestyle{fancy}
\fancyhf{} % Clear all headers/footers
\lhead{Assignment 2}
\rfoot{\thepage}
\renewcommand{\headrulewidth}{0.4pt}
\renewcommand{\footrulewidth}{0pt}


% --- Title and TOC ---

\begin{document}

\title{4P78 Assignment 2}
\author{
    Parker TenBroeck 7376726\\
    pt21zs@brocku.ca
}
\date{\today}

\makeatletter
\begin{titlepage}
	\def \LOGOPATH {brock.jpg}
	\def \UNIVERSITY {Brock University}
	\def \FACULTY {Faculty of Mathematics \& Science}
	\def \DEPARTMENT {Department of Computer Science}
	\def \COURSETITLE {COSC 3P93: Parallel Computing}
	\def \SUPERVISOR {Robson De Grande}
	
	
	\vfill
	\begin{center}
		\includegraphics[width=0.6\textwidth]{brock.jpg}
		\fontsize{14pt}{14pt}\selectfont
		\vfill
		\UNIVERSITY \\
		\FACULTY \\
		\DEPARTMENT \\
		\vfill
		\fontsize{18pt}{18pt}\selectfont
		\textbf{\COURSETITLE} \\[0.5cm]
		\textbf{\@title}
		\vfill
		\fontsize{14pt}{14pt}\selectfont
		Prepared By: \\[0.5cm]
		
		\begin{tabular}[t]{c}
			\@author
		\end{tabular}\par
	
	    \vfill
		Instructor: \\
		\SUPERVISOR
		\vfill
		\@date
	\end{center}
\end{titlepage}
\makeatother

\newpage


\section{Performance}

\subsection{}

\begin{enumerate}
	\item Define Task Latency and determine the Average Task Latency for the entire batch of 10 tasks \\
	how long it takes for a system to start or complete a task after the task is issued.\\
	$$
	\frac{0+0+0+0+1+2+2.5+3+3.5+5}{10}=1.7s
	$$
	
	\item Define Task Throughput and determine the Overall Task Throughput in the observed time interval \\
	
	The number of tasks the system can complete in a givin duration of time.\\
	$$
	\frac{10}{7}=1.43s
	$$
\end{enumerate}

\section{Parallel Design Patterns}

\subsection{}
\begin{enumerate}
	\item 
	I3 : A[10] = X + 10 (T1 write) → I4 : Y = A[10] + 5 (T2 read)\\
	True (Flow) a write to A[10] is later read by the other thread, so the flow of data is from I3 to I4
	
	\item I1 : P = P + 1 (T1 write) and I2 : P = P + 2 (T2 write) \\
	Output (Write–Write) both instructions write the same scalar P\\~\\
	I3 : X = A[5] (T1 read) → I4 : A[5] = 2 × Y (T2 write)\\
	(Read→Write) — I3 reads A[5] and a later instruction writes A[5], so the read precedes the write (anti).\\
	
	\item I3 : M[0][1] = 1 (T1 write) and I4 : M[0][1] = 2 (T2 write)\\
	Output (Write–Write) — both write the same matrix element, so there is an output dependency.
	
	\item No Dependency's for any instructions
	
	\item I3 : A[8] = Y (T1 write) → I4 : Y = A[8] + 1 (T2 read of A[8])\\
	Type: True (Flow) — I3 writes A[8] and I4 reads that location.\\~\\
	
	I3 : A[8] = Y (T1 reads Y) and I4 : Y = A[8] + 1 (T2 writes Y)\\
	Anti-Dependency (Read→Write) on Y — I3 reads Y while I4 later writes Y.
\end{enumerate}

\subsection{}
\begin{enumerate}
	\item 
\end{enumerate}

\subsection{}
\begin{enumerate}
	\item 
\end{enumerate}

\subsection{}
\begin{enumerate}
	\item 
\end{enumerate}

\section{OpenMPI}

\subsection{}
\begin{enumerate}
	\item $[6]$ MPI\_Init: starts the MPI environment
	\item $[7]$ MPI\_Comm\_rank:  gives each process its rank (ID) in MPI\_COMM\_WORLD
	\item $[8]$ MPI\_Comm\_size: gives the total number of processes
	\item $[16]$ MPI\_Sendrecv: Sends left\_send (1 double) to left\_neighbor with tag 0.\\
	Simultaneously receives right\_recv (1 double) from right\_neighbor with tag 0.\\
	Sendrecv is useful in situations where two processes both want to exchange data but it is not clear who should go first.
	\item $[20]$ MPI\_Sendrecv: Sends right\_send to right\_neighbor with tag 1.\\
	Simultaneously Receives left\_recv from left\_neighbor with tag 1.
	
	\item $[27]$ MPI\_Barrier: is a global synchronization, all processes must reach the boundary before any can proceed past it
	
	\item $[31]$ MPI\_Finalize: Cleans up the MPI environment
\end{enumerate}

\subsection{}
\begin{enumerate}
	\item $[9]$ MPI\_Init: starts the MPI environment
	\item $[10]$ MPI\_Comm\_rank:  gives each process its rank (ID) in MPI\_COMM\_WORLD
	\item $[11]$ MPI\_Comm\_size: gives the total number of processes
	\item $[16]$ MPI\_Scatter: distribute the dataset by splitting global\_data into size blocks of length local\_N and sends a block to each process
	\item $[21]$ MPI\_Bcast: Broadcast from rank 0 to all other processes so each process has a copy of the centroid array.
	\item $[25]$ MPI\_Reduce: each process gives its local\_sum and the reduce operation sums all the values together.
	\item $[32]$ MPI\_Gather: collect all local means on rank 0
	\item $[36]$ MPI\_Allreduce: Reduce all values with MAX operation (so find the max of all values) and share it to all ranks
	\item $[40]$ MPI\_Allgather: Collect all the sums from every process into the array.
	
	\item $[44]$ MPI\_Finalize: Cleans up the MPI environment
\end{enumerate}

\section{}

\section{Tools}
\subsection{}
\begin{enumerate}
	\item Factors for sequential performance include, time complecity, quality of implementation, CPU speed, ILP, Memory behavior(cache, memory bandwidth), Compiler optimizations, Instruction throughput
	\item Factors for parallel performance include, Parallel fraction of the algorithm, load balancing, synchronization overhead, communication overhead, memory contention, scalability, thread mamagement overhead, cache coherence overhead.
\end{enumerate}
\section{Parallel Program Design}

\subsection{}
The methodological design of parallel programs is a structured approach that transforms a problem into an efficient parallel solution by identifying parallel work, organizing tasks, and minimizing communication and synchronization so the program scales and performs well on multiple processors\\
\begin{enumerate}
	\item Partitioning
	\item Communication
	\item Agglomeration
	\item Mapping
	\item Synchronization
	\item Performance Tuning
\end{enumerate}


%\nocite{*}
%\bibliographystyle{IEEEtran}
%\bibliography{references}

\end{document}
